% =====================================================================
%  IEEE Transactions on Electron Devices -- REVISED MANUSCRIPT
%  TED-2026-08-2206-R  (resubmission)
%
%  MARKED copy  : \newcommand{\rev}[1]{{\color{blue}#1}}
%  CLEAN copy   : \newcommand{\rev}[1]{#1}
%  Compile: pdflatex -> bibtex -> pdflatex -> pdflatex
% =====================================================================
\documentclass[journal]{IEEEtran}
\usepackage{graphicx}
\usepackage{amsmath,amssymb}
\usepackage[T1]{fontenc}
\usepackage{url}
\usepackage{booktabs}
\usepackage{siunitx}
\usepackage{xcolor}
\sisetup{detect-all}

%\newcommand{\rev}[1]{{\color{blue}#1}}     % MARKED
\newcommand{\rev}[1]{#1}                    % CLEAN

\newcommand{\hzo}{Hf$_{0.5}$Zr$_{0.5}$O$_2$}
\newcommand{\Prem}{P_r}
\newcommand{\Ec}{E_c}
\newcommand{\Vc}{V_c}
\newcommand{\Vth}{V_{\mathrm{th}}}
\newcommand{\Vimp}{V_{\mathrm{imp}}}
\newcommand{\Neq}{N_{\mathrm{eq}}}
\newcommand{\dVh}{\Delta V_{\mathrm{hyst}}}
\newcommand{\Cil}{C_{\mathrm{IL}}}
\newcommand{\Cfe}{C_{\mathrm{HZO}}}
\newcommand{\Cacc}{C_{\mathrm{acc}}}
\newcommand{\Rs}{R_s}
\newcommand{\ucm}{\micro\coulomb\per\centi\meter\squared}
\newcommand{\uFcm}{\micro\farad\per\centi\meter\squared}
\newcommand{\MVcm}{\mega\volt\per\centi\meter}
\newcommand{\figph}[1]{\fbox{\parbox[c][3cm][c]{0.95\linewidth}
  {\centering\ttfamily\detokenize{#1}\\[2pt]\rmfamily\itshape not found}}}
\newcommand{\figmaybe}[2]{\IfFileExists{#1}{\includegraphics[width=#2]{#1}}{\figph{#1}}}

\begin{document}

\title{\rev{Sign Reversal of the DC-Sweep Capacitance--Voltage Hysteresis
Window in \hzo{} Ferroelectric Transistors}}

\author{Apu~Das, Agniva~Paul, Debraj~Barman, Zhao-Feng~Lou,
        Yannick~Raffel, Min-Hung~Lee, and~Sourav~De
\thanks{This work was partly funded by the National Science and Technology
Council (NSTC), Taiwan, under grants 115-2221-E-007-144, by the Taiwan
Semiconductor Research Institute (TSRI) grant JDP114-Y1-002, and by
PSMC-115A0330.}
\thanks{Apu Das, Agniva Paul, Debraj Barman, and Sourav De are with the
College of Semiconductor Research, National Tsing Hua University, Hsinchu
30044, Taiwan (e-mail: apu.das@gapp.nthu.edu.tw; sourav.de@mx.nthu.edu.tw).}
\thanks{Zhao-Feng Lou and Min-Hung Lee are with National Taiwan University,
Taipei 10617, Taiwan. Yannick Raffel is with the Center Nanoelectronic
Technologies, Fraunhofer IPMS, Dresden 01109, Germany.}}

\markboth{IEEE Transactions on Electron Devices}%
{Das \MakeLowercase{\textit{et al.}}: Sign Reversal of the DC-Sweep C--V Hysteresis Window}
\maketitle

% =====================================================================
\begin{abstract}
\rev{Memory-window collapse in HfO$_2$ ferroelectric transistors is
commonly attributed to ferroelectric fatigue. We show that the DC-sweep
$C$--$V$ hysteresis window of TaN/\hzo{}(\SI{10}{\nano\meter})/
SiO$_2$(\SI{0.85}{\nano\meter})/p-Si transistors does not merely close
but \emph{reverses sign}, becoming clockwise beyond
$5\times10^2$--$2\times10^3$ cycles. The reversal is reproduced on five
independent devices, is unchanged when the measurement integration time
is increased ten-fold, and is present at every frequency from
\SI{10}{\kilo\hertz} to \SI{1}{\mega\hertz} and every sweep amplitude
from \SI{\pm1.5}{\volt} to \SI{\pm3}{\volt}. Transfer characteristics
confirm it: $\Delta\Vth$ falls from \SI{+0.581}{\volt} to
\SI{-0.015}{\volt} at $10^4$ cycles, independently of the extraction
criterion. Over the same range the remanent polarization
\emph{increases}, rising 45\% by $10^2$ cycles and remaining above its
pristine value at $10^4$. The window efficiency consequently falls from
$\sim$16\% to $-4$\%. Interfacial charge first compensates, then
over-compensates, the polarization field at the channel.}
\end{abstract}

\begin{IEEEkeywords}
Ferroelectric field-effect transistor, hafnium oxide, endurance,
hysteresis window, charge trapping, capacitance--voltage, interfacial
layer, PUND.
\end{IEEEkeywords}

% =====================================================================
\section{Introduction}
\IEEEPARstart{F}{erroelectricity} in doped HfO$_2$ \cite{Fe1,Fe2} made a
CMOS-compatible ferroelectric field-effect transistor (FeFET) practical,
and integrated FeFET technologies now exist at the \SI{28}{\nano\meter}
node and below \cite{Khan2020,QLC_FeNAND}. Endurance remains the
principal obstacle: the memory window typically collapses between $10^4$
and $10^5$ write cycles \cite{Mulaosmanovic_2021,HKG_FeFET,HfO_SCM},
and the gap persists across silicon- and zirconium-doped stacks alike
\cite{endurance1,Device_life}.

The default explanation is ferroelectric fatigue \cite{fatigue}. An
alternative attributes the collapse to charge injected into the
ferroelectric/interfacial-layer region, which screens the polarization
from the channel without destroying it
\cite{Ref_DD_1,Charge_Trap_endurance,Toprasertpong,passlack}.
\rev{Charge trapping, trap generation and interfacial-layer effects are
established features of FeFET operation rather than mechanisms
identified here \cite{Tasneem,Trapping_dynamics,Defect_dynamics}.} The
two mechanisms imply opposite remedies --- ferroelectric materials
engineering \cite{Stability} versus interface engineering
\cite{endurance3,endurance4,AgnivaTSA} --- so distinguishing them
matters. Reported evidence supports both, including oxygen-vacancy
redistribution during cycling \cite{Oxygen_Vac,ref_1_OV,Ref_2OV}.

\rev{Two of our own recent studies bear directly on this work and we
state the relationship explicitly. In \cite{ACSAMI} we showed that
ferroelectric polarization survives after the FeFET memory window
collapses. In \cite{ACSAELM} we extracted an apparent channel-coupled
charge density from the threshold-voltage shift, defined a
polarization-to-channel coupling factor by comparison with the switching
charge, and reported its decrease from 3.5\% to 0.9\% at $10^4$ cycles.
Both establish that the coupling between polarization and channel
degrades while polarization persists; we do not claim that observation
as new. A companion preprint reports the stability of the gate-stack
polarization response in degraded devices \cite{preprint}.}

\rev{Three findings in the present work are not contained in those
studies. First, the hysteresis window does not merely approach zero --- it
becomes \emph{negative}. A coupling factor is non-negative by
construction, so the regime in which interfacial charge
\emph{over-compensates} rather than merely screens the polarization
field cannot be expressed in that framework. Beyond the crossing, a
positive gate sweep leaves the transistor harder to turn on, the
opposite of its pristine behavior. Second, the reversal is established
on five independent devices, across three decades of measurement
frequency, four sweep amplitudes and a ten-fold change in integration
time. Third, it is resolved from the gate-stack $C$--$V$ alone, and
independently confirmed by transfer characteristics.}

% =====================================================================
\section{Device Fabrication and Measurement}

\subsection{Gate Stack}
The device is a p-Si channel FeFET with a
TaN/\hzo{}(\SI{10}{\nano\meter})/SiO$_2$(\SI{0.85}{\nano\meter}) gate
stack [Fig.~\ref{fig:stack}(a)]. The chemical SiO$_2$ interfacial layer
(IL) is unavoidable on a silicon channel and is the region of interest.
The insulator stack forms a series capacitive divider,
\begin{equation}
\frac{V_{\mathrm{HZO}}}{V_{\mathrm{ins}}} = \frac{\Cil}{\Cil+\Cfe},
\label{eq:divider}
\end{equation}
where $V_{\mathrm{ins}}=V_{\mathrm{HZO}}+V_{\mathrm{IL}}$. With
$\varepsilon_{\mathrm{HZO}}=30$ and $\varepsilon_{\mathrm{IL}}=3.9$,
$\Cfe=\SI{2.66}{\uFcm}$ and $\Cil=\SI{4.06}{\uFcm}$, giving 0.60.
Because the semiconductor capacitance is in series and neglected in
\eqref{eq:divider}, 0.60 is an upper bound.

\rev{\subsection{Devices and Data Provenance}
Five nominally identical devices from the same die and process are
reported, denoted device-1 to device-3 and device-6, device-7. Device-1
is the device of the original submission; the others were not previously
measured. All were cycled under identical conditions.

Two measurement types were performed on single devices and we state this
explicitly. Polarization (PUND) was measured on one device of the die;
because polarization behavior is uniform across devices of this process
\cite{ACSAMI,ACSAELM}, the resulting $2\Prem(N)$ is used as a
\emph{common reference} against which the window of every device is
compared, rather than as a per-device quantity. Transfer characteristics
were measured on a single device. Neither is presented as a
device-to-device reproducibility result; that role is carried by the
$C$--$V$ data on five devices.}

\subsection{Polarization Measurement}
Switching-current and $P$--$E$ measurements used triangular waveforms
with a half-period of \SI{500}{\micro\second}. Polarization was
evaluated by positive-up-negative-down (PUND) measurement:
\begin{equation}
Q_{\mathrm{sw}}^{+}=\int\!\left[I_P-I_U\right]dt,\quad
Q_{\mathrm{sw}}^{-}=\int\!\left[I_N-I_D\right]dt,
\label{eq:pund}
\end{equation}
where $P$, $N$ denote switching and $U$, $D$ nonswitching pulses, and
\begin{equation}
2\Prem=\frac{|Q_{\mathrm{sw}}^{+}|+|Q_{\mathrm{sw}}^{-}|}{2A}.
\label{eq:2pr}
\end{equation}
Coercive voltages were taken from the $P$--$V$ zero crossings, and
\begin{equation}
\Vimp=\frac{\Vc^{+}+\Vc^{-}}{2},\qquad
\Ec=\frac{|\Vc^{+}|+|\Vc^{-}|}{2d}.
\label{eq:imprint}
\end{equation}
For the pristine device $\Vc=\SI{2.08}{\volt}$, which with
\eqref{eq:divider} gives $\Ec\approx\SI{1.25}{\MVcm}$, consistent with
reported values \cite{Ferro_switching,imprint3}.

\subsection{Cycling, Capacitance and Transfer Measurement}
Cycling used \SI{\pm5}{\volt}, \SI{50}{\micro\second} pulses. $C$--$V$
was measured on a Keysight B1500A CMU in $C_p$--$R_p$ mode with a
\SI{20}{\milli\volt} AC level. \rev{Devices 6 and 7 were characterized at
seven frequencies (\SI{1}{\kilo\hertz}--\SI{1}{\mega\hertz}) and four
sweep amplitudes (\SI{\pm1.5}{\volt}--\SI{\pm3}{\volt}) at each
checkpoint; devices 2 and 3 at \SI{100}{\kilo\hertz} and
\SI{\pm3}{\volt} with two integration settings.} Source, drain and body
were tied so that the inversion branch is supplied with minority
carriers.

Two features of the capacitance measurement require explicit statement.
The sweep begins at the negative rail and integrates for one power-line
cycle per point, so a full sweep takes $\sim$\SI{1}{\second} --- four
orders of magnitude slower than the write pulse. Consequently the pulsed
state does not survive to mid-sweep, and $C$--$V$ curves acquired after
program and after erase are indistinguishable. The reported quantity is
therefore the \emph{DC-sweep hysteresis window}
\begin{equation}
\dVh = V_{\mathrm{fwd}} - V_{\mathrm{rev}}
       \Big|_{C=C_{\min}+\frac{1}{2}(C_{\max}-C_{\min})}.
\label{eq:window}
\end{equation}
\rev{It is not the retained state of a pulse-written bit, and no claim
about stored information is made from it alone; the corresponding
transistor quantity is measured separately in Section~III-E.} The sign
convention is fixed throughout: $\dVh>0$ denotes a counter-clockwise
loop (the ferroelectric sense); $\dVh<0$ a clockwise loop (the
electron-trapping sense).

\rev{Transfer characteristics were measured with an SPGU write
(\SI{\pm5}{\volt}, \SI{50}{\micro\second}) and an SMU read
($V_D=\SI{0.1}{\volt}$, $V_G$ swept $-1.5$ to \SI{+2.5}{\volt}).
$\Vth$ was extracted at a constant current of \SI{1}{\micro\ampere}.
Because the gate must be switched between the pulse and read
configurations through a switching matrix, the read delay was
\SI{60}{\second}, determined from measurement timestamps; all terminals
were held at \SI{0}{\volt} during that interval and each read was
preceded by a fresh write.}

\rev{\subsection{Integration-Time Control}
To test whether the reported hysteresis is affected by charge leaking
during the ADC integration window, the cycling series was repeated on
devices 2 and 3 with the integration coefficient increased from 1 to 10
power-line cycles, i.e.\ \SI{16.7}{\milli\second} to
\SI{167}{\milli\second} per point, all else unchanged.}

\subsection{Series-Resistance Correction and Screening}
A series resistance was extracted from the Nicollian--Brews
strong-accumulation expression \cite{nicollian1982,schroder2006}
\begin{equation}
\Rs = \frac{G_{ma}}{G_{ma}^{2}+\omega^{2}C_{ma}^{2}},
\label{eq:rs}
\end{equation}
and capacitances corrected using
\begin{equation}
C_c=\frac{(G^{2}+\omega^{2}C^{2})C}{a^{2}+\omega^{2}C^{2}},\quad
a=G-(G^{2}+\omega^{2}C^{2})\Rs .
\label{eq:corr}
\end{equation}
Sweeps with median loss tangent $D=1/\omega R_pC_p$ above the stated
screen were excluded from window statistics.

\subsection{Figures of Merit}
An unscreened remanent polarization would produce a window
$2\Prem/\Cil$. We define the window efficiency
\begin{equation}
\eta = \frac{\dVh}{2\Prem/\Cil},
\label{eq:eta}
\end{equation}
\rev{using the common reference $2\Prem(N)$ of Section~II-B. Where a
window deficit is attributed to interfacial charge, an \emph{equivalent}
areal density follows from}
\begin{equation}
\rev{\Delta\Neq=\frac{\Cil\,\Delta V}{q}
 \approx\Delta V\times\SI{2.5e13}{\per\centi\meter\squared\per\volt}.}
\label{eq:neq}
\end{equation}
\rev{Equation \eqref{eq:neq} assumes a sheet charge at the interfacial
layer and an unchanged $\Cil$; it is an equivalent density inferred from
a voltage shift, not a directly extracted interface-trap density, and is
denoted $\Neq$ for that reason.}

% =====================================================================
\section{Results and Discussion}

\subsection{The Ferroelectric Remains Intact to $10^4$ Cycles}
Fig.~\ref{fig:fe} shows the polarization response. Field cycling
produces a pronounced wake-up, consistent with oxygen-vacancy
redistribution and domain-wall de-pinning
\cite{cho2024,Trapping_dynamic2}: $2\Prem$ rises from \SI{20.2}{\ucm}
pristine to \SI{29.2}{\ucm} at $10^2$ cycles, a 45\% increase, and the
positive switching peak nearly doubles. Both remain at or above their
pristine values through $10^4$. $\Vc$ is constant to within 4\% at every
verified checkpoint through $10^4$ (Table~\ref{tab:switch}).

\rev{Wake-up and charge trapping proceed simultaneously rather than
sequentially: devices 2 and 3 show wake-up in the hysteresis window
itself, rising from \SI{0.822}{\volt} to \SI{0.920}{\volt} and from
\SI{0.804}{\volt} to \SI{0.968}{\volt} between pristine and one cycle
before collapsing.}

Degradation appears only at $10^5$ cycles, where $2\Prem$ falls to
\SI{18.6}{\ucm}, the peaks broaden, and the non-switching current rises.
$\Vc$ \emph{decreases} there, from \SI{2.08}{\volt} to
\SI{1.94}{\volt} --- not a domain-pinning signature --- consistent with a
leakier, electrically softened film \cite{Poole_Frankel_2}.
\rev{We therefore do not claim that ferroelectric fatigue is absent. The
ferroelectric does degrade, approximately two decades after the
hysteresis window has collapsed, so loss of switchable polarization
alone cannot account for the early collapse.}

The imprint voltage drifts monotonically from \SI{-48}{\milli\volt}
pristine to \SI{+9}{\milli\volt} at $10^4$, crossing zero between $10^3$
and $10^4$, consistent with accumulation of a negative charge sheet on
the channel side of the ferroelectric \cite{imprint3}.

\subsection{The Window Collapses and Reverses Sign}
Fig.~\ref{fig:cvfam} shows the $C$--$V$ evolution. The separation
between forward and reverse branches narrows monotonically, closes near
$10^3$ cycles, and reopens at $10^4$ with the branches exchanged.

\rev{Fig.~\ref{fig:window}(a) shows $\dVh$ for five independent devices.
All exhibit a counter-clockwise window when fresh and a clockwise window
by $10^4$ cycles; interpolated zero crossings are 503, 871, 1117, 2136
and 2422 cycles [Fig.~\ref{fig:window}(b)], a spread of a factor of
five. At $10^4$ the window is $-0.15$ to $-\SI{0.28}{\volt}$ across all
five (Table~\ref{tab:window}). Device-1 has a smaller pristine window
(\SI{0.363}{\volt}) than the other four (0.79--\SI{0.82}{\volt}),
reflecting die-to-die variation; the sign reversal is nonetheless common
to all.}

\rev{The window measured on the inversion flank --- the flank
corresponding to $\Vth$ --- collapses and inverts on the same schedule
[Fig.~\ref{fig:window}(c)], excluding an artifact of the extraction
point, and is consistently larger than the flatband-flank window,
indicating that the loop stretches rather than translating rigidly. A
rigid shift, as produced by a fixed charge sheet, would displace both
flanks equally. The reversal is also present at every sweep amplitude
[Fig.~\ref{fig:window}(d)]; we note that the 50\%-of-range extraction
criterion is not directly comparable between amplitudes and therefore
report only the persistence of the sign change, not a trend in
magnitude.}

\rev{\subsection{The Reversal Is Not a Measurement Artifact}
Three independent checks are shown in Fig.~\ref{fig:artifact}. First,
$\dVh$ is suppressed at low frequency and saturates above
$\sim$\SI{100}{\kilo\hertz}, following}
\begin{equation}
\dVh(f) = \Delta V_\infty - \frac{\delta V}{1+(f/f_c)^{2}},
\label{eq:debye}
\end{equation}
\rev{with $f_c=\SI{25.6}{\kilo\hertz}$ (device-6) and
\SI{16.8}{\kilo\hertz} (device-7) in the pristine state, both consistent
with the 22--\SI{34}{\kilo\hertz} obtained on device-1
[Fig.~\ref{fig:artifact}(b)]. Fits at $10^2$ cycles use only four points
on a partly collapsed window and are shown as open symbols; we draw no
conclusion from them. The reversal itself is present at all seven
measurement frequencies from \SI{1}{\kilo\hertz} to \SI{1}{\mega\hertz}.

Second, increasing the ADC integration ten-fold changes $\dVh$ by
15--22\% but does not change its sign at any checkpoint through $10^5$
on either device [Fig.~\ref{fig:artifact}(c)]. In the reversed regime
the magnitude moves \emph{toward} zero at longer integration. Charge
leaking during the integration window would produce the opposite
dependence, since a longer integration would deepen an artifact-driven
negative window. The observed trend is therefore inconsistent with that
mechanism.

Third, the reversal is present in uncorrected $C_p$ data: of the sweeps
at $10^4$ cycles on device-1 for which the flank crossing is defined
without any series correction, 25 of 27 are negative, spanning five
frequencies, four amplitudes and both pulse states.}

\subsection{Gate-Stack Controls}
\rev{Fig.~\ref{fig:controls} shows the accumulation capacitance, loss
spectra and conduction for devices 6 and 7. The validity of the
correction of \eqref{eq:corr} is checkpoint-dependent, and we restrict
the claim accordingly: a physically meaningful corrected capacitance is
obtained through $10^3$ cycles on device-1 and through $10^2$ on
device-2, while at $10^4$ cycles $\omega\Rs C\approx0.50$ and the
correction is ill-conditioned. No corrected value is quoted there, and
the claim of an unchanged capacitance is not extended to the checkpoint
at which the sign reversal occurs.

Within the valid range the result is nonetheless decisive. On device-2,
over $N=1$ to $10^2$ cycles, $\Cacc=184.2\pm\SI{1.6}{\pico\farad}$
($\pm0.9$\%) with the correction below 1.9\% at every point and
$D=0.12$--0.15, while $\dVh$ falls from \SI{0.920}{\volt} to
\SI{0.331}{\volt} --- 64\% of the collapse. The stack capacitance is
therefore constant to better than 1\%, on essentially uncorrected data,
over the range in which the window loses most of its magnitude.}

\subsection{Transistor-Level Confirmation}
\rev{Fig.~\ref{fig:transistor} shows transfer characteristics measured
after program and erase at each checkpoint. Defining
$\Delta\Vth=\Vth(\text{erase})-\Vth(\text{program})$ so that the sign
convention matches $\dVh$, the transistor window falls from
\SI{+0.581}{\volt} pristine to \SI{+0.243}{\volt} at $10^2$,
\SI{+0.108}{\volt} at $10^3$, and \SI{-0.015}{\volt} at $10^4$
[Fig.~\ref{fig:transistor}(b)]. The reversal is independent of the
extraction criterion, being present at 1, 3, 10, 30 and
\SI{100}{\micro\ampere} [Fig.~\ref{fig:transistor}(c)]. Both threshold
states converge and cross, falling from 0.29/\SI{0.87}{\volt} pristine
to 0.083/\SI{0.067}{\volt} at $10^4$.

The relation between $\dVh$ and the pulse-written threshold window,
stated as an assumption in Section~II-D, is therefore supported by
direct measurement: the two quantities collapse on the same schedule and
change sign at the same cycle count [Fig.~\ref{fig:decouple}(c)]. We
note that the two are measured on very different timescales ---
\SI{60}{\second} for the transistor read against
$\sim$\SI{17}{\milli\second} per point for $C$--$V$ --- so their
agreement indicates that the responsible charge is stable over at least
that range. The off-state current floor of
$\sim$\SI{3e-7}{\ampere} in these devices prevents extraction of the
subthreshold swing.}

\subsection{Polarization and Window Are Decoupled}
\rev{Fig.~\ref{fig:decouple} places the measurements together and
constitutes the central result. Normalising each device's window to its
own pristine value [Fig.~\ref{fig:decouple}(a)], all five collapse to
between $-0.26$ and $-0.42$ of their initial magnitude at $10^4$ cycles,
while the reference $2\Prem$ rises to 1.45 of its pristine value at
$10^2$ and remains above unity at $10^4$.

The window efficiency of \eqref{eq:eta}, evaluated against the common
reference polarization, falls from 15.9--16.5\% pristine to $-3.2$ to
$-4.2$\% at $10^4$ cycles for devices 2, 3, 6 and 7
[Fig.~\ref{fig:decouple}(b), Table~\ref{tab:eta}]; device-1 follows the
same trajectory from a lower starting value. The agreement among four
independent devices at both ends of the range is closer than the
device-to-device spread in the pristine window itself.}

\subsection{Interpretation}
\rev{At $10^4$ cycles the polarization is present and larger than
pristine, and both the capacitance and transistor windows have reversed
sign. Over the range where the series correction is valid, the
gate-stack capacitance is unchanged to better than 1\% while most of the
collapse occurs. These observations are consistent with charge
accumulating within the channel-side gate stack: such charge can screen
the polarization from the channel while leaving both the polarization
and the series capacitance intact.} From \eqref{eq:neq}, the
\SI{0.51}{\volt} swing in $\dVh$ on device-1 corresponds to
$\Delta\Neq\approx\SI{1.3e13}{\per\centi\meter\squared}$, a density
readily accommodated at an HfO$_2$/SiO$_2$/Si interface
\cite{Ref_DD_1,Toprasertpong}, and comparable to trap densities
extracted by charge pumping and low-frequency noise on nominally
identical stacks \cite{Trapping_dynamics,Defect_dynamics}.

The sign reversal further shows that this charge does not merely
neutralize the polarization's influence. It exceeds it, and there is
necessarily an intermediate cycle count at which the two contributions
cancel exactly.

\rev{The present measurements do not determine where the compensating
charge resides. They do not distinguish charge in bulk HZO, at the
HZO/SiO$_2$ interface, within SiO$_2$, or at the SiO$_2$/Si interface,
and we describe it only as effective charge within the channel-side gate
stack. Nor do we claim that the flank asymmetry, imprint drift and
conduction increase share a single origin. The mechanism is presented as
consistent with charge trapping rather than as an identification of the
responsible traps; transmission electron microscopy with electron
energy-loss spectroscopy, or charge pumping across the cycling range,
would be required to localize and quantify it directly.}

\rev{Three limitations remain. The five devices are from one die, and
the zero-crossing cycle count varies by a factor of five among them.
Polarization and transfer characteristics were each measured on a single
device (Section~II-B), so the decoupling argument combines a
five-device window measurement with single-device polarization and
transistor references. And $\eta$ depends on the convention chosen for
the ideal window in \eqref{eq:eta}; alternative denominators shift the
absolute values but not the sign change or the trend.}

% =====================================================================
\section{Conclusion}
\rev{The DC-sweep $C$--$V$ hysteresis window of an HfO$_2$ FeFET does
not merely close with cycling: it reverses sign. The reversal is
reproduced on five independent devices with zero crossings between
$5\times10^2$ and $2\times10^3$ cycles, survives a ten-fold increase in
measurement integration time, is present at all seven measurement
frequencies and four sweep amplitudes, and is confirmed by transfer
characteristics in which $\Delta\Vth$ changes sign at the same cycle
count. Over the same range the remanent polarization stands above its
pristine value, and where the lumped $C_p$--$R_p$ model is valid the
gate-stack capacitance is constant to better than 1\% while 64\% of the
collapse occurs. Loss of switchable polarization alone therefore cannot
account for the early collapse of the hysteresis window in these
devices, and interfacial charge that first compensates and then
over-compensates the polarization field provides a consistent account.}

\section*{Acknowledgment}
The authors thank the Taiwan Semiconductor Research Institute and the
Nano-Fabrication Facility, National Yang-Ming Chiao Tung University, and
Prof.~Min-Hung Lee for the measurement facility at NTU.

\bibliographystyle{IEEEtran}
\bibliography{references}

% =====================================================================
\begin{figure}[!t]\centering
\figmaybe{FIG1_stack.pdf}{\columnwidth}
\caption{Device and measurement protocol. (a)~Gate stack and equivalent
series divider, \eqref{eq:divider}; at most 60\% of the gate bias
reaches the ferroelectric, and with $\Vc=\SI{2.08}{\volt}$ this gives
$\Ec\approx\SI{1.25}{\MVcm}$. (b)~At each checkpoint the device is
cycled with \SI{\pm5}{\volt}, \SI{50}{\micro\second} pulses, then
characterized. The $C$--$V$ sweep starts at the negative rail and takes
$\sim$\SI{17}{\milli\second} per point, four orders of magnitude slower
than the write pulse; the reported window is DC-sweep hysteresis, not a
pulse-written state.}
\label{fig:stack}\end{figure}

\begin{figure}[!t]\centering
\figmaybe{FIG2_ferroelectric.pdf}{\columnwidth}
\caption{Ferroelectric response versus cycling. (a)~$P$--$V$ loops and
(b)~switching current. (c)~$2\Prem$ from \eqref{eq:2pr} and switching
peaks; dotted line marks the pristine $2\Prem$. (d)~$\Vimp$ and $\Vc$
from \eqref{eq:imprint}. Wake-up raises $2\Prem$ 45\% by $10^2$ cycles
and it stays above pristine through $10^4$, with $\Vc$ constant to 4\%.
At $10^5$, $\Vc$ decreases as the peaks broaden --- not a domain-pinning
signature. \rev{Measured on a single device of the die; used as the
common polarization reference (Section~II-B).} Open symbol at $10^3$:
single unverified trace.}
\label{fig:fe}\end{figure}

\begin{figure}[!t]\centering
\figmaybe{FIG3_CV_families.pdf}{\columnwidth}
\caption{\rev{$C$--$V$ evolution of device-7 at \SI{500}{\kilo\hertz},
\SI{\pm3}{\volt}. Solid: forward sweep; dashed: reverse. The
forward--reverse separation narrows monotonically, closes near $10^3$
cycles, and reopens at $10^4$ with the branches exchanged.}}
\label{fig:cvfam}\end{figure}

\begin{figure}[!t]\centering
\figmaybe{FIG4_window.pdf}{\columnwidth}
\caption{\rev{The window collapses and reverses sign. (a)~$\dVh$ versus
cycling for five independent devices; shaded region marks clockwise
loops. Device-1 (grey) has a smaller pristine window but the same
trajectory. (b)~Interpolated zero crossings: 503, 871, 1117, 2136 and
2422 cycles. (c)~Flatband and threshold flanks, devices 6 and 7; both
invert on the same schedule. (d)~Window versus sweep amplitude at
\SI{1}{\mega\hertz}; the reversal is present at every amplitude.}}
\label{fig:window}\end{figure}

\begin{figure}[!t]\centering
\figmaybe{FIG5_artifact.pdf}{\columnwidth}
\caption{\rev{The reversal is not a measurement artifact.
(a)~Inverted-Debye dispersion, \eqref{eq:debye}, devices 6 and 7.
(b)~Fitted $f_c$; filled symbols are fits with $\geq5$ points, open
symbols the four-point fits at $10^2$ cycles which are not reliable.
Grey band: device-1 range. (c)~PLC 1 versus PLC 10
(\SI{16.7} versus \SI{167}{\milli\second} per point). The sign is
unchanged at every checkpoint through $10^5$; magnitudes move toward
zero at longer integration, opposite to the dependence expected from a
leakage artifact.}}
\label{fig:artifact}\end{figure}

\begin{figure}[!t]\centering
\figmaybe{FIG6_controls.pdf}{\columnwidth}
\caption{\rev{Gate-stack controls, devices 6 and 7.
(a)~Accumulation capacitance at \SI{500}{\kilo\hertz}. (b)~Loss spectra;
dotted line marks the screening threshold. (c)~Accumulation conductance
versus cycling.}}
\label{fig:controls}\end{figure}

\begin{figure}[!t]\centering
\figmaybe{FIG7_decoupling.pdf}{\columnwidth}
\caption{\rev{Polarization and window are decoupled. (a)~$\dVh$
normalised to each device's own pristine value, five devices, with
$2P_{r}/2P_{r,0}$ from the common reference measurement.
(b)~Window efficiency \eqref{eq:eta} evaluated against that reference:
15.9--16.5\% pristine falling to $-3.2$ to $-4.2$\% at $10^4$ for
devices 2, 3, 6 and 7. (c)~Normalised $\Delta\Vth$ (transistor, single
device, \SI{60}{\second} read delay) plotted against the normalised
$\dVh$ of all five devices.}}
\label{fig:decouple}\end{figure}

\begin{figure}[!t]\centering
\figmaybe{FIG8_transistor.pdf}{\columnwidth}
\caption{\rev{Transistor-level confirmation, single device.
(a)~Transfer characteristics after program and erase, pristine and
$10^4$ cycles, \SI{60}{\second} read delay; dotted line marks the
\SI{1}{\micro\ampere} extraction criterion. (b)~$\Vth$ of both states,
which converge and cross. (c)~$\Delta\Vth$ at five current criteria; the
sign reversal is criterion-independent. (d)~$\Delta\Vth$ and $\dVh$
versus cycling.}}
\label{fig:transistor}\end{figure}

% =====================================================================
\begin{table}[!t]
\caption{\rev{DC-Sweep Hysteresis Window Across Five Devices}}
\label{tab:window}\centering\footnotesize\setlength{\tabcolsep}{4pt}
\begin{tabular}{lccccc}
\toprule
$N$ & dev-1 & dev-2 & dev-3 & dev-6 & dev-7\\
\midrule
pristine & $+0.363$ & $+0.822$ & $+0.804$ & $+0.791$ & $+0.803$\\
$1$      & $+0.265$ & $+0.920$ & $+0.968$ & $+0.719$ & $+0.740$\\
$10$     & $+0.221$ & $+0.806$ & ---      & $+0.606$ & $+0.672$\\
$10^2$   & $+0.122$ & $+0.331$ & ---      & $+0.396$ & $+0.492$\\
$10^3$   & $-0.052$ & $-0.021$ & $+0.129$ & $+0.010$ & $+0.172$\\
$10^4$   & $-0.151$ & $-0.265$ & $-0.262$ & $-0.205$ & $-0.275$\\
$10^5$   & ---      & $-0.115$ & $-0.105$ & ---      & $-0.116$\\
\midrule
zero cross.\ & 503 & 871 & 2136 & 1117 & 2422\\
\bottomrule
\end{tabular}\\[2pt]
\parbox{\columnwidth}{\scriptsize All values in volts, from
\eqref{eq:window}. Devices 1--3 at \SI{100}{\kilo\hertz}, devices 6--7
at \SI{500}{\kilo\hertz}, all at \SI{\pm3}{\volt}. Entries omitted where
the loss tangent exceeds the screening threshold.}
\end{table}

\begin{table}[!t]
\caption{\rev{Window Efficiency $\eta$ (\%) Against the Common
Polarization Reference}}
\label{tab:eta}\centering\footnotesize\setlength{\tabcolsep}{4pt}
\begin{tabular}{lcccccc}
\toprule
$N$ & $2\Prem$ & dev-1 & dev-2 & dev-3 & dev-6 & dev-7\\
    & (\si{\ucm}) & & & & &\\
\midrule
pristine & 20.2 & 7.3 & 16.5 & 16.2 & 15.9 & 16.1\\
$1$      & 25.1 & 4.3 & 14.9 & 15.7 & 11.6 & 12.0\\
$10$     & 27.7 & 3.2 & 11.8 & ---  & 8.9  & 9.9\\
$10^2$   & 29.2 & 1.7 & 4.6  & ---  & 5.5  & 6.8\\
$10^3$   & 22.1 & $-1.0$ & $-0.4$ & 2.4 & 0.2 & 3.2\\
$10^4$   & 26.5 & $-2.3$ & $-4.1$ & $-4.0$ & $-3.2$ & $-4.2$\\
$10^5$   & 18.6 & ---    & $-2.5$ & $-2.3$ & ---    & $-2.5$\\
\bottomrule
\end{tabular}\\[2pt]
\parbox{\columnwidth}{\scriptsize $\eta$ from \eqref{eq:eta}. The
$2\Prem$ column is the common reference measured on a single device of
the die (Section~II-B) and is applied to all five devices; $\eta$
therefore compares each device's window against a representative stack
polarization rather than its own.}
\end{table}

\begin{table}[!t]
\caption{Ferroelectric Switching Diagnostics}
\label{tab:switch}\centering\footnotesize\setlength{\tabcolsep}{4pt}
\begin{tabular}{lcccccc}
\toprule
$N$ & $\Vc$ & $\Ec$ & $\Vimp$ & $I_{sw}^{+}$ & $I_{sw}^{-}$ & $I_{ns}$\\
    & (V) & (\si{\MVcm}) & (mV) & (\si{\micro\ampere})
    & (\si{\micro\ampere}) & (\si{\micro\ampere})\\
\midrule
pristine & 2.08 & 1.25 & $-48$ & 28.2 & $-32.4$ & 3.15\\
$1$      & 2.07 & 1.24 & $-48$ & 45.5 & $-59.3$ & 1.46\\
$10$     & 2.10 & 1.26 & $-28$ & 49.4 & $-64.1$ & 0.83\\
$10^2$   & 2.16 & 1.30 & $-23$ & 51.3 & $-65.9$ & 0.62\\
$10^3$   & 2.37 & 1.42 & $-12$ & 33.9 & $-40.7$ & 1.52\\
$10^4$   & 2.07 & 1.24 & $+9$  & 54.0 & $-56.1$ & 1.20\\
$10^5$   & 1.94 & 1.16 & $-35$ & 41.2 & $-38.6$ & 1.87\\
\bottomrule
\end{tabular}\\[2pt]
\parbox{\columnwidth}{\scriptsize From \eqref{eq:imprint};
$\Ec=0.60\,\Vc/t_{\mathrm{HZO}}$ using \eqref{eq:divider}. The $10^3$
row derives from a single unverified trace and is excluded from the
quoted $\Vc$ variation.}
\end{table}

\begin{table}[!t]
\caption{\rev{Transistor Threshold Window, Single Device}}
\label{tab:vth}\centering\footnotesize\setlength{\tabcolsep}{5pt}
\begin{tabular}{lcccc}
\toprule
$N$ & $\Vth$(prog.) & $\Vth$(erase) & $\Delta\Vth$ & $I_{on}/I_{off}$\\
    & (V) & (V) & (V) & \\
\midrule
pristine & 0.294 & 0.875 & $+0.581$ & 1053\\
$1$      & 0.421 & 0.939 & $+0.518$ & 1050\\
$10$     & 0.550 & 0.923 & $+0.373$ & 1063\\
$10^2$   & 0.533 & 0.776 & $+0.243$ & 1025\\
$10^3$   & 0.327 & 0.434 & $+0.108$ & 969\\
$10^4$   & 0.083 & 0.067 & $-0.015$ & 538\\
\bottomrule
\end{tabular}\\[2pt]
\parbox{\columnwidth}{\scriptsize $\Vth$ at a constant current of
\SI{1}{\micro\ampere}; $\Delta\Vth=\Vth(\text{erase})-\Vth(\text{program})$
so that the sign convention matches $\dVh$. Read delay \SI{60}{\second}.}
\end{table}

\end{document}
